\documentclass[../cours_1re.tex]{subfiles}
\begin{document}
\setcounter{chapter}{13}
\chapter{Radian et cercle trigonométrique}
\label{ch:14}
\begin{objectif}
Définir les radians et les fonctions sinus et cosinus sur tous les réels ; démontrer les valeurs remarquables et les relations d'angles associés. Prérequis : géométrie du cercle et Pythagore.
\end{objectif}
\section{Mesurer un angle par un arc}
Sur un cercle de rayon $r>0$, un angle au centre interceptant un arc de longueur $s$ a pour mesure en radians $s/r$.
Ce quotient ne dépend pas du rayon : une homothétie multiplie l'arc et le rayon par le même facteur.
Un tour mesure $2\pi$ radians, un demi-tour $\pi$, un quart de tour $\pi/2$.
La conversion est $\theta_{\mathrm{rad}}=\theta_{\mathrm{deg}}\pi/180$.
Pour un angle non orienté $\theta\ge0$, $s=r\theta$ ; pour un déplacement orienté, le signe indique le sens et la longueur parcourue vaut $r|\theta|$.
\section{Enrouler la droite sur le cercle}
Dans un repère orthonormé, le cercle trigonométrique a centre $O$ et rayon $1$. On part de $I(1;0)$ ; les angles positifs tournent dans le sens contraire des aiguilles d'une montre.
Le réel $t$ correspond au point atteint après une longueur orientée $t$. Les réels $t+2k\pi$, $k\in\Z$, donnent le même point.
\begin{definition}
Le point associé à $t$ a pour coordonnées $(\cos t;\sin t)$.
Ainsi $\cos^2t+\sin^2t=1$, et les deux coordonnées appartiennent à $[-1;1]$.
\end{definition}
L'identité est l'équation du cercle unité, issue de Pythagore. Dans le premier quadrant, elle retrouve les rapports cosinus et sinus d'un triangle rectangle d'hypoténuse $1$.
\begin{center}\begin{tikzpicture}[scale=2]
\draw[->](-1.3,0)--(1.4,0) node[right]{$x$};
\draw[->](0,-1.3)--(0,1.4) node[above]{$y$};
\draw[accent,thick](0,0)circle(1);
\coordinate(M) at (40:1);
\draw[thick](0,0)--(M) node[above right]{$M(\cos t;\sin t)$};
\draw[dashed](M)--({cos(40)},0) node[below]{$\cos t$};
\draw[dashed](M)--(0,{sin(40)}) node[left]{$\sin t$};
\draw[->](.3,0)arc(0:40:.3) node[right]{$t$};
\end{tikzpicture}\end{center}
\section{Valeurs remarquables démontrées}
À $0,\pi/2,\pi,3\pi/2$, lire les intersections avec les axes.
Pour $\pi/4$, le triangle rectangle est isocèle : ses deux côtés de l'angle droit ont même longueur $a$ et $2a^2=1$. Comme $a>0$, $a=\sqrt2/2$.
Pour $\pi/3$, prendre un triangle équilatéral de côté $1$. La hauteur le coupe en deux triangles rectangles, avec un côté $1/2$ et une hauteur $\sqrt{1-1/4}=\sqrt3/2$. Cela donne $\cos(\pi/3)=1/2$ et $\sin(\pi/3)=\sqrt3/2$. L'angle complémentaire donne les valeurs en $\pi/6$.
\begin{center}\begin{tabular}{c|ccccc}
$t$&$0$&$\pi/6$&$\pi/4$&$\pi/3$&$\pi/2$\\\hline
$\cos t$&$1$&$\sqrt3/2$&$\sqrt2/2$&$1/2$&$0$\\
$\sin t$&$0$&$1/2$&$\sqrt2/2$&$\sqrt3/2$&$1$
\end{tabular}\end{center}
\section{Symétries et angles associés}
\begin{propriete}
\[
\begin{aligned}
\cos(-t)&=\cos t,&\sin(-t)&=-\sin t,\\
\cos(\pi-t)&=-\cos t,&\sin(\pi-t)&=\sin t,\\
\cos(t+\pi)&=-\cos t,&\sin(t+\pi)&=-\sin t,\\
\cos(\pi/2-t)&=\sin t,&\sin(\pi/2-t)&=\cos t.
\end{aligned}
\]
Sinus et cosinus sont $2\pi$-périodiques.
\end{propriete}
\begin{demonstration}
La réflexion dans l'axe horizontal change $(x;y)$ en $(x;-y)$ ; dans l'axe vertical, en $(-x;y)$ ; le demi-tour change les deux signes. La réflexion dans la droite $y=x$ échange les coordonnées. Ces transformations correspondent respectivement aux angles $-t,\pi-t,t+\pi,\pi/2-t$. Un tour complet ne change pas le point.
\end{demonstration}
\begin{exemple}
$\cos(5\pi/6)=-\sqrt3/2$, $\sin(5\pi/6)=1/2$.
Pour $-13\pi/4$, ajouter $4\pi$ donne $3\pi/4$ : cosinus $-\sqrt2/2$, sinus $\sqrt2/2$.
Réduire modulo $2\pi$ signifie ici enlever un nombre entier de tours.
\end{exemple}
\section{Équations sur le cercle, en complément}
Deux points du cercle ont même abscisse s'ils sont égaux ou symétriques par rapport à l'axe horizontal. Donc
$\cos x=\cos a$ équivaut à $x=2k\pi+a$ ou $x=2k\pi-a$.
De même $\sin x=\sin a$ équivaut à $x=a+2k\pi$ ou $x=\pi-a+2k\pi$.
On impose ensuite l'intervalle de résolution. Par exemple, sur $[0;2\pi[$,
$\sin x=1/2$ donne $x=\pi/6$ ou $5\pi/6$.
\begin{finchapitre}
\section{Approcher pi par des polygones}
\approfondissement Pour un polygone régulier inscrit à $n\ge3$ côtés de longueur $s$ dans le cercle unité,
\[
\frac{ns}{2}<\pi<\frac{ns}{2\sqrt{1-s^2/4}}.
\]
La borne inférieure est la moitié du périmètre inscrit ; la supérieure est la moitié du périmètre du polygone circonscrit. Le triangle rectangle de demi-côté $s/2$ donne l'apothème $\sqrt{1-s^2/4}$ ; des triangles semblables donnent le demi-côté circonscrit $(s/2)/\sqrt{1-s^2/4}$.
La comparaison des périmètres utilise les propriétés usuelles de longueur du cercle.
Après doublement du nombre de côtés, Pythagore donne
\[
s_{\mathrm{nouveau}}=\sqrt{2-\sqrt{4-s^2}}
=\frac{s}{\sqrt{2+\sqrt{4-s^2}}}.
\]
La seconde écriture évite la soustraction de deux nombres très proches. On démarre à $n=6,s=1$, d'où $3<\pi<2\sqrt3$, puis on resserre l'encadrement. Le fichier Python réalise ce calcul.


\section{Utilité et prolongements}
Les radians rendent les formules de longueur et les futures dérivées trigonométriques cohérentes. Le cercle donne simultanément valeurs, signes, symétries et solutions.

En licence, sinus et cosinus décrivent les rotations et les oscillations ; en master, les séries de Fourier décomposent des signaux en modes trigonométriques. Le passage d'un phénomène périodique à un angle sur un cercle prépare déjà cette représentation.
\begin{retenir}
Vérifier le mode radians de la calculatrice. Ramener un angle à un tour, repérer son quadrant, puis utiliser les symétries ; ne pas deviner les signes.
\end{retenir}
\end{finchapitre}
\end{document}
